\appendix
\section{Anexos: scripts y configuración}
\label{anx:devsecops-config}

Copia de los archivos del repositorio \url{https://github.com/jeffreymonjacastro/MisiOps} (rama \texttt{develop}) citados en el informe.

\subsection{Workflow \texttt{.github/workflows/devsecops.yml}}
\label{anx:devsecops}

\begin{lstlisting}
name: devsecops

# Security pipeline (specs/011-devsecops). Gates block pull requests into main and pushes to main;
# into develop they only report (constitution exemption EX-001). Jobs and gates:
# specs/011-devsecops/contracts/pipeline-jobs.md

on:
  pull_request:
    branches: [develop, main]
  push:
    branches: [develop, main]
  workflow_dispatch:

permissions:
  contents: read

concurrency:
  group: devsecops-${{ github.ref }}
  cancel-in-progress: true

env:
  ENFORCE: ${{ github.base_ref == 'main' || github.ref == 'refs/heads/main' }}
  SEMGREP_IMAGE: semgrep/semgrep:1.183.0
  GITLEAKS_IMAGE: zricethezav/gitleaks:v8.30.1
  TRUFFLEHOG_IMAGE: trufflesecurity/trufflehog:3.99.8
  TRIVY_IMAGE: aquasec/trivy:0.76.1
  ZAP_IMAGE: ghcr.io/zaproxy/zaproxy:2.17.0

jobs:
  build:
    runs-on: ubuntu-24.04-arm  # linux/arm64 images for Fargate Graviton (specs/012-aws-deployment R2)
    steps:
      - uses: actions/checkout@v4
      - name: Build images once (scanned, attacked and delivered as the same bytes)
        run: |
          docker build -t misiops-backend:latest backend
          docker build -t misiops-frontend:latest frontend
          docker save misiops-backend:latest -o backend.tar
          docker save misiops-frontend:latest -o frontend.tar
      - uses: actions/upload-artifact@v4
        with:
          name: images
          path: "*.tar"
          retention-days: 1

  sast:
    runs-on: ubuntu-latest
    steps:
      - uses: actions/checkout@v4
      - name: Semgrep (backend + frontend source)
        run: |
          mkdir -p reports
          docker run --rm -u "$(id -u):$(id -g)" -e HOME=/tmp -v "$PWD:/src" -w /src "$SEMGREP_IMAGE" \
            semgrep scan --config p/default --config p/python --config p/typescript --metrics=off \
            --json-output reports/semgrep.json --text-output reports/semgrep.txt backend frontend
      - uses: actions/upload-artifact@v4
        if: always()
        with:
          name: report-sast
          path: reports/
          retention-days: 7
      - name: Gate (Semgrep ERROR)
        run: security/scripts/gate.sh Semgrep "$(jq '[.results[] | select(.extra.severity == "ERROR")] | length' reports/semgrep.json)"

  secrets:
    runs-on: ubuntu-latest
    steps:
      - uses: actions/checkout@v4
        with:
          fetch-depth: 0
      - name: Gitleaks (full history)
        run: |
          mkdir -p reports
          docker run --rm -u "$(id -u):$(id -g)" -v "$PWD:/repo" -w /repo "$GITLEAKS_IMAGE" \
            git --config .gitleaks.toml --redact --exit-code 0 \
            --report-format json --report-path reports/gitleaks.json .
      - name: TruffleHog (full history)
        run: |
          docker run --rm -u "$(id -u):$(id -g)" -e HOME=/tmp -v "$PWD:/repo" "$TRUFFLEHOG_IMAGE" \
            git file:///repo --json --no-update --exclude-paths=/repo/security/trufflehog-exclude.txt > reports/trufflehog.json
      - uses: actions/upload-artifact@v4
        if: always()
        with:
          name: report-secrets
          path: reports/
          retention-days: 7
      - name: Gate (any secret)
        run: |
          security/scripts/gate.sh Gitleaks "$(jq length reports/gitleaks.json)"
          security/scripts/gate.sh TruffleHog "$(grep -c '"SourceMetadata"' reports/trufflehog.json || true)"

  sca:
    runs-on: ubuntu-latest
    steps:
      - uses: actions/checkout@v4
      - name: Trivy DB cache key
        id: day
        run: echo "date=$(date -u +%F)" >> "$GITHUB_OUTPUT"
      - uses: actions/cache@v4
        with:
          path: ~/.cache/trivy
          key: trivy-db-${{ steps.day.outputs.date }}
          restore-keys: trivy-db-
      - name: Trivy dependencies (uv.lock, pnpm-lock.yaml) + SBOM
        run: |
          mkdir -p reports ~/.cache/trivy
          trivy() { docker run --rm -u "$(id -u):$(id -g)" -e TRIVY_CACHE_DIR=/cache -v "$HOME/.cache/trivy:/cache" -v "$PWD:/src" -w /src "$TRIVY_IMAGE" "$@"; }
          # docs/ holds archived reports (incl. an SBOM), not project dependencies.
          trivy fs --scanners vuln --skip-dirs docs --ignorefile .trivyignore.yaml --format json --output reports/trivy-deps.json .
          trivy fs --scanners vuln --skip-dirs docs --ignorefile .trivyignore.yaml --skip-db-update --format table --output reports/trivy-deps.txt .
          trivy fs --skip-dirs docs --skip-db-update --format cyclonedx --output reports/sbom.cdx.json .
      - uses: actions/upload-artifact@v4
        if: always()
        with:
          name: report-sca
          path: reports/
          retention-days: 7
      - name: Gate (fixable HIGH/CRITICAL)
        run: security/scripts/gate.sh "Trivy dependencies" "$(jq -f security/scripts/trivy-blocking.jq reports/trivy-deps.json)"

  container:
    needs: build
    runs-on: ubuntu-latest
    steps:
      - uses: actions/checkout@v4
      - uses: actions/download-artifact@v4
        with:
          name: images
      - name: Trivy DB cache key
        id: day
        run: echo "date=$(date -u +%F)" >> "$GITHUB_OUTPUT"
      - uses: actions/cache@v4
        with:
          path: ~/.cache/trivy
          key: trivy-db-${{ steps.day.outputs.date }}
          restore-keys: trivy-db-
      - name: Trivy images (backend, frontend, postgres)
        run: |
          mkdir -p reports ~/.cache/trivy
          trivy() { docker run --rm -u "$(id -u):$(id -g)" -e TRIVY_CACHE_DIR=/cache -v "$HOME/.cache/trivy:/cache" -v "$PWD:/src" -w /src "$TRIVY_IMAGE" "$@"; }
          postgres=$(grep -oE 'image: postgres[^[:space:]]*' docker/docker-compose.yaml | cut -d' ' -f2)
          trivy image --ignorefile .trivyignore.yaml --input backend.tar --format json --output reports/trivy-image-backend.json
          trivy image --ignorefile .trivyignore.yaml --skip-db-update --input frontend.tar --format json --output reports/trivy-image-frontend.json
          trivy image --ignorefile .trivyignore.yaml --skip-db-update --format json --output reports/trivy-image-postgres.json "$postgres"
      - uses: actions/upload-artifact@v4
        if: always()
        with:
          name: report-container
          path: reports/
          retention-days: 7
      - name: Gate (fixable HIGH/CRITICAL, secrets)
        run: |
          rc=0
          for image in backend frontend postgres; do
            security/scripts/gate.sh "Trivy image $image" "$(jq -f security/scripts/trivy-blocking.jq "reports/trivy-image-$image.json")" || rc=1
          done
          exit $rc

  iac:
    runs-on: ubuntu-latest
    steps:
      - uses: actions/checkout@v4
      - name: Trivy config (Dockerfiles) + Semgrep (docker-compose)
        run: |
          mkdir -p reports
          docker run --rm -u "$(id -u):$(id -g)" -e TRIVY_CACHE_DIR=/tmp/trivy -v "$PWD:/src" -w /src "$TRIVY_IMAGE" \
            config --skip-dirs docs --format json --output reports/trivy-config.json .
          # Trivy has no Compose scanner, so the orchestration file goes through Semgrep's compose rules.
          docker run --rm -u "$(id -u):$(id -g)" -e HOME=/tmp -v "$PWD:/src" -w /src "$SEMGREP_IMAGE" \
            semgrep scan --config p/docker-compose --metrics=off --json-output reports/semgrep-compose.json docker
      - uses: actions/upload-artifact@v4
        if: always()
        with:
          name: report-iac
          path: reports/
          retention-days: 7
      - name: Gate (HIGH/CRITICAL misconfigurations)
        run: |
          rc=0
          security/scripts/gate.sh "Trivy config" "$(jq -f security/scripts/trivy-blocking.jq reports/trivy-config.json)" || rc=1
          security/scripts/gate.sh "Semgrep compose" "$(jq '[.results[] | select(.extra.severity == "ERROR")] | length' reports/semgrep-compose.json)" || rc=1
          exit $rc

  dast:
    needs: build
    runs-on: ubuntu-24.04-arm  # linux/arm64 images for Fargate Graviton (specs/012-aws-deployment R2)
    steps:
      - uses: actions/checkout@v4
      - uses: actions/download-artifact@v4
        with:
          name: images
      - name: Start the stack from the scanned images
        run: |
          docker load -i backend.tar && docker load -i frontend.tar
          # Throwaway credentials generated per run; never committed (FR-018).
          pw=$(openssl rand -hex 24)
          printf 'POSTGRES_PASSWORD=%s\nSECRET_KEY=%s\n' "$pw" "$(openssl rand -hex 32)" > docker/.env.dev
          printf '%s' "$pw" > docker/pg_password.txt
          docker compose -f docker/docker-compose.yaml up -d --no-build
          curl -fsS --retry 30 --retry-all-errors --retry-delay 2 -o /dev/null http://localhost:8000/
          curl -fsS --retry 30 --retry-all-errors --retry-delay 2 -o /dev/null http://localhost:3000/
      - name: Get a bearer token for the authenticated API scan
        run: |
          email="zap-${GITHUB_RUN_ID}@example.com" pass=$(openssl rand -hex 16)
          body=$(jq -nc --arg e "$email" --arg p "$pass" '{name: "ZAP scanner", email: $e, password: $p}')
          curl -fsS -X POST http://localhost:8000/api/v1/auth/register -H 'Content-Type: application/json' -d "$body" -o /dev/null
          token=$(curl -fsS -X POST http://localhost:8000/api/v1/auth/login -H 'Content-Type: application/json' -d "$body" | jq -r .access_token)
          echo "::add-mask::$token"
          echo "ZAP_TOKEN=$token" >> "$GITHUB_ENV"
      - name: ZAP baseline (frontend) and API scan (OpenAPI, authenticated)
        run: |
          mkdir -p reports && chmod 777 reports
          cp security/zap/rules.conf reports/zap-rules.conf
          zap() { docker run --rm --network host -v "$PWD/reports:/zap/wrk:rw" "$@"; }
          zap "$ZAP_IMAGE" zap-baseline.py -t http://localhost:3000 -c zap-rules.conf -I \
            -x zap-frontend.xml -r zap-frontend.html -J zap-frontend.json || true
          zap -e ZAP_AUTH_HEADER_VALUE="Bearer $ZAP_TOKEN" -e ZAP_AUTH_HEADER_SITE=localhost:8000 "$ZAP_IMAGE" \
            zap-api-scan.py -t http://localhost:8000/openapi.json -f openapi -c zap-rules.conf -I \
            -x zap-api.xml -r zap-api.html -J zap-api.json -z "-config scanner.maxScanDurationInMins=10" || true
          # A missing report means ZAP never reached the app: fail instead of reporting "no findings".
          test -s reports/zap-frontend.json && test -s reports/zap-api.json
      - uses: actions/upload-artifact@v4
        if: always()
        with:
          name: report-dast
          path: reports/
          retention-days: 7
      - name: Gate (High-risk alerts)
        run: |
          rc=0
          for scan in frontend api; do
            security/scripts/gate.sh "ZAP $scan" "$(jq '[.site[]?.alerts[]? | select(.riskcode == "3")] | length' "reports/zap-$scan.json")" || rc=1
          done
          exit $rc
      - name: Stop the stack
        if: always()
        run: docker compose -f docker/docker-compose.yaml down -v

  scripts-test:
    runs-on: ubuntu-latest
    steps:
      - uses: actions/checkout@v4
      - run: bash security/tests/test_upload_to_defectdojo.sh

  summary:
    needs: [sast, secrets, sca, container, iac, dast]
    if: always()
    runs-on: ubuntu-latest
    steps:
      - uses: actions/download-artifact@v4
        with:
          pattern: report-*
          path: reports
          merge-multiple: true
      - name: Findings by tool and severity
        run: |
          cd reports
          sev() { jq -r "$1" "$2" 2>/dev/null | sort | uniq -c | awk '{printf "%s %s, ", $1, $2}'; }
          trivy='.Results[]? | (.Vulnerabilities[]?, .Misconfigurations[]?) | .Severity'
          {
            echo "## Security findings (mode: $([ "$ENFORCE" = true ] && echo blocking || echo report-only))"
            echo "| Report | Findings by severity |"
            echo "| --- | --- |"
            for f in semgrep.json semgrep-compose.json; do echo "| $f | $(sev '.results[].extra.severity' "$f") |"; done
            echo "| gitleaks.json | $(jq length gitleaks.json 2>/dev/null) secrets |"
            echo "| trufflehog.json | $(grep -c '"SourceMetadata"' trufflehog.json 2>/dev/null) secrets |"
            for f in trivy-*.json; do echo "| $f | $(sev "$trivy" "$f") |"; done
            for f in zap-frontend.json zap-api.json; do echo "| $f | $(sev '.site[]?.alerts[]? | .riskdesc | split(" ")[0]' "$f") |"; done
          } >> "$GITHUB_STEP_SUMMARY"

  defectdojo:
    needs: [sast, secrets, sca, container, iac, dast]
    if: always()
    continue-on-error: true
    runs-on: ubuntu-latest
    env:
      DEFECTDOJO_URL: ${{ secrets.DEFECTDOJO_URL }}
      DEFECTDOJO_TOKEN: ${{ secrets.DEFECTDOJO_TOKEN }}
    steps:
      - uses: actions/checkout@v4
      - uses: actions/download-artifact@v4
        with:
          pattern: report-*
          path: reports
      - name: Import every report (reimport-scan)
        run: |
          if [ -z "$DEFECTDOJO_URL" ]; then
            echo "skipped: no DefectDojo configured (use security/scripts/import_run.sh ${GITHUB_RUN_ID})"
            exit 0
          fi
          security/scripts/import_run.sh reports

  deliver:
    needs: [sast, secrets, sca, container, iac, dast, scripts-test]
    if: github.event_name == 'push' && github.ref == 'refs/heads/main'
    runs-on: ubuntu-latest
    permissions:
      contents: read
      packages: write
    steps:
      - uses: actions/download-artifact@v4
        with:
          name: images
      - name: Publish the scanned images to GHCR
        env:
          GH_TOKEN: ${{ secrets.GITHUB_TOKEN }}
        run: |
          echo "$GH_TOKEN" | docker login ghcr.io -u "${{ github.actor }}" --password-stdin
          for app in backend frontend; do
            docker load -i "$app.tar"
            for tag in "$GITHUB_SHA" latest; do
              docker tag "misiops-$app:latest" "ghcr.io/jeffreymonjacastro/misiops-$app:$tag"
              docker push "ghcr.io/jeffreymonjacastro/misiops-$app:$tag"
            done
          done
\end{lstlisting}

\subsection{Gate por herramienta: \texttt{security/scripts/gate.sh}}
\label{anx:gate}

\begin{lstlisting}
#!/usr/bin/env bash
# Security gate (FR-012). Blocks only when ENFORCE=true (pull requests into main, pushes to main);
# otherwise warns, because merges into develop are report-only under constitution exemption EX-001.
# Usage: gate.sh <tool> <blocking-finding-count>
set -euo pipefail

if [ $# -ne 2 ]; then
  echo "usage: $0 <tool> <count>" >&2
  exit 2
fi
tool=$1 count=$2
mode=report-only
[ "${ENFORCE:-false}" = true ] && mode=blocking
[ -n "${GITHUB_STEP_SUMMARY:-}" ] && echo "- **$tool**: $count blocking findings ($mode)" >> "$GITHUB_STEP_SUMMARY"

if [ "$count" -eq 0 ]; then
  echo "$tool: no blocking findings"
elif [ "$mode" = blocking ]; then
  echo "::error::$tool: $count blocking findings"
  exit 1
else
  echo "::warning::$tool: $count findings would block a merge into main (report-only, constitution EX-001)"
fi
\end{lstlisting}

\subsection{Filtro de Trivy: \texttt{security/scripts/trivy-blocking.jq}}
\label{anx:jq}

\begin{lstlisting}
# Blocking findings in a Trivy JSON report (FR-012): fixable HIGH/CRITICAL vulnerabilities,
# HIGH/CRITICAL misconfigurations and any secret. Unfixed vulnerabilities are reported only.
[
  .Results[]?
  | (.Vulnerabilities[]? | select((.Severity == "HIGH" or .Severity == "CRITICAL") and ((.FixedVersion // "") != ""))),
    (.Misconfigurations[]? | select(.Severity == "HIGH" or .Severity == "CRITICAL")),
    (.Secrets[]?)
] | length
\end{lstlisting}

\subsection{Subida a DefectDojo: \texttt{security/scripts/upload\_to\_defectdojo.sh}}
\label{anx:upload}

\begin{lstlisting}
#!/usr/bin/env bash
# Upload one scan report to DefectDojo through /api/v2/reimport-scan/.
# Contract: specs/011-devsecops/contracts/upload-script.md
# Usage: upload_to_defectdojo.sh <scan_type> <test_title> <report_file>
set -euo pipefail

if [ $# -lt 3 ]; then
  echo "usage: $0 <scan_type> <test_title> <report_file>" >&2
  exit 2
fi
scan_type=$1 test_title=$2 report=$3

# A missing file means the scanner did not run; a zero-finding report is still uploaded
# so close_old_findings can close what an earlier run found.
if [ ! -f "$report" ]; then
  echo "skip: $report not found"
  exit 0
fi
for var in DEFECTDOJO_URL DEFECTDOJO_TOKEN; do
  if [ -z "${!var:-}" ]; then
    echo "missing variable: $var" >&2
    exit 1
  fi
done

tmp=$(mktemp -d)
trap 'rm -rf "$tmp"' EXIT

# TruffleHog writes log lines next to findings; DefectDojo only parses the findings.
if [ "$scan_type" = "Trufflehog Scan" ]; then
  grep '"SourceMetadata"' "$report" > "$tmp/trufflehog.json" || true
  report=$tmp/trufflehog.json
fi

# Header from a file keeps the token off the process list and out of logs.
printf 'Authorization: Token %s\n' "$DEFECTDOJO_TOKEN" > "$tmp/auth"

form=(
  -F "scan_type=$scan_type"
  -F "test_title=$test_title"
  -F "file=@$report"
  -F "product_name=${DEFECTDOJO_PRODUCT:-MisiOps}"
  -F "product_type_name=${DEFECTDOJO_PRODUCT_TYPE:-Aplicaciones Web}"
  -F "engagement_name=${DEFECTDOJO_ENGAGEMENT:-CI - ${GITHUB_HEAD_REF:-${GITHUB_REF_NAME:-local}}}"
  -F "auto_create_context=true"
  -F "close_old_findings=true"
  -F "active=true"
  -F "verified=false"
)
[ -n "${GITHUB_RUN_ID:-}" ] && form+=(-F "build_id=$GITHUB_RUN_ID")
[ -n "${GITHUB_SHA:-}" ] && form+=(-F "commit_hash=$GITHUB_SHA")
[ -n "${GITHUB_REF_NAME:-}" ] && form+=(-F "branch_tag=${GITHUB_HEAD_REF:-$GITHUB_REF_NAME}")

status=$(curl -sS -o "$tmp/body" -w '%{http_code}' -X POST \
  "${DEFECTDOJO_URL%/}/api/v2/reimport-scan/" -H "@$tmp/auth" "${form[@]}")

if [ "$status" = "201" ]; then
  echo "ok: $test_title"
else
  echo "error $status: $(cat "$tmp/body")" >&2
  exit 1
fi
\end{lstlisting}

\subsection{Mapa de reportes y parsers: \texttt{security/scripts/reports.tsv}}
\label{anx:reports}

\begin{lstlisting}
# scan_type    test_title    report file (FR-010, distinct titles: research R7)
Semgrep JSON Report    Semgrep - source    semgrep.json
Gitleaks Scan    Gitleaks - history    gitleaks.json
Trufflehog Scan    TruffleHog - history    trufflehog.json
Trivy Scan    Trivy - dependencies    trivy-deps.json
Trivy Scan    Trivy - image backend    trivy-image-backend.json
Trivy Scan    Trivy - image frontend    trivy-image-frontend.json
Trivy Scan    Trivy - image postgres    trivy-image-postgres.json
Trivy Scan    Trivy - config    trivy-config.json
ZAP Scan    ZAP - frontend baseline    zap-frontend.xml
ZAP Scan    ZAP - API    zap-api.xml
Semgrep JSON Report    Semgrep - compose    semgrep-compose.json
\end{lstlisting}

\subsection{Hooks locales: \texttt{.pre-commit-config.yaml}}
\label{anx:precommit}

\begin{lstlisting}
# Local secret blocking (FR-009). Run by .githooks/pre-commit (core.hooksPath); needs pre-commit 4.7.0 or uv.
repos:
  - repo: https://github.com/gitleaks/gitleaks
    rev: v8.30.1
    hooks:
      - id: gitleaks
  - repo: https://github.com/pre-commit/pre-commit-hooks
    rev: v6.0.0
    hooks:
      - id: detect-private-key
\end{lstlisting}

\subsection{Reglas de ZAP: \texttt{security/zap/rules.conf}}
\label{anx:zap}

\begin{lstlisting}
# ZAP rule overrides: <rule id><TAB><IGNORE|WARN|FAIL><TAB><name>
# The pipeline gate counts High-risk alerts from the JSON report (FR-012), so this file
# only holds suppressions. Every IGNORE line needs a comment above it:
# id: <rule> | justification: <why> | review: <YYYY-MM-DD>   (FR-012b)
\end{lstlisting}

\subsection{Workflow \texttt{.github/workflows/infra-ci.yml}}
\label{anx:infraci}

\begin{lstlisting}
name: infra-ci

# Infrastructure gate (specs/012-aws-deployment, FR-007, NFR-007): tests and an offline synth of the CDK app.
# No AWS credentials: the stacks have no context lookups and the region is fixed in infra/bin/misiops.ts.

on:
  pull_request:
    branches: [develop]
    paths: ["infra/**", ".github/workflows/infra-ci.yml"]
  push:
    branches: [develop]
    paths: ["infra/**", ".github/workflows/infra-ci.yml"]

permissions:
  contents: read

jobs:
  infra-ci:
    runs-on: ubuntu-latest
    timeout-minutes: 5
    defaults:
      run:
        working-directory: infra
    steps:
      - uses: actions/checkout@v4
      - uses: pnpm/action-setup@v4
        with:
          version: 9.6.0
      - uses: actions/setup-node@v4
        with:
          node-version: 20
          cache: pnpm
          cache-dependency-path: infra/pnpm-lock.yaml
      - run: pnpm install --frozen-lockfile
      - name: Unit tests (stack assertions + smoke script)
        run: pnpm test
      - name: Synth dev and prod without credentials
        run: |
          pnpm synth -c env=dev -c imageTag=ci -c alertEmail=ci@example.com --quiet
          pnpm synth -c env=prod -c imageTag=ci -c alertEmail=ci@example.com --quiet
          test "$(grep -c '"SecretString"' cdk.out/MisiOps-dev.template.json)" = 0
\end{lstlisting}

\subsection{Workflow \texttt{.github/workflows/deploy.yml}}
\label{anx:deploy}

\begin{lstlisting}
name: deploy

# Deploys the images scanned by the devsecops run that triggered this workflow (specs/012-aws-deployment).
# develop -> dev (automatic, after the pipeline completes; gates are report-only, EX-001)
# main    -> prod (waits for the required reviewer of the GitHub Environment "prod")
# Contract: specs/012-aws-deployment/contracts/deploy-workflow.md

on:
  workflow_run:
    workflows: [devsecops]
    types: [completed]
    branches: [develop, main]

permissions:
  contents: read
  actions: read
  id-token: write

env:
  AWS_REGION: us-east-1

jobs:
  resolve:
    runs-on: ubuntu-latest
    outputs:
      env: ${{ steps.map.outputs.env }}
      deploy: ${{ steps.map.outputs.deploy }}
      sha: ${{ github.event.workflow_run.head_sha }}
      run_id: ${{ github.event.workflow_run.id }}
    steps:
      - id: map
        env:
          BRANCH: ${{ github.event.workflow_run.head_branch }}
          CONCLUSION: ${{ github.event.workflow_run.conclusion }}
          EVENT: ${{ github.event.workflow_run.event }}
        run: |
          case "$BRANCH" in
            develop) env=dev ;;
            main) env=prod ;;
            *) env= ;;
          esac
          deploy=false
          if [ -n "$env" ] && [ "$CONCLUSION" = success ] && [ "$EVENT" = push ]; then deploy=true; fi
          echo "env=$env" >> "$GITHUB_OUTPUT"
          echo "deploy=$deploy" >> "$GITHUB_OUTPUT"
          echo "branch=$BRANCH conclusion=$CONCLUSION event=$EVENT -> env=${env:-none} deploy=$deploy" | tee -a "$GITHUB_STEP_SUMMARY"

  deploy:
    needs: resolve
    if: needs.resolve.outputs.deploy == 'true'
    runs-on: ubuntu-latest
    timeout-minutes: 30
    environment: ${{ needs.resolve.outputs.env }}
    concurrency:
      group: deploy-${{ needs.resolve.outputs.env }}
      cancel-in-progress: false
    env:
      ENV_NAME: ${{ needs.resolve.outputs.env }}
      SHA: ${{ needs.resolve.outputs.sha }}
    steps:
      - name: Start timer
        run: echo "START=$(date +%s)" >> "$GITHUB_ENV"

      - uses: actions/checkout@v4
        with:
          ref: ${{ needs.resolve.outputs.sha }}

      - name: Download the scanned images of the triggering devsecops run (FR-009)
        uses: actions/download-artifact@v8
        with:
          name: images
          run-id: ${{ needs.resolve.outputs.run_id }}
          github-token: ${{ secrets.GITHUB_TOKEN }}
          path: images

      - name: Assume the deploy role through OIDC (FR-011)
        uses: aws-actions/configure-aws-credentials@v6.3.0
        with:
          role-to-assume: ${{ vars.AWS_DEPLOY_ROLE_ARN }}
          aws-region: ${{ env.AWS_REGION }}

      - id: ecr
        uses: aws-actions/amazon-ecr-login@v2

      - name: Push the images to ECR tagged with the commit SHA
        id: push
        env:
          REGISTRY: ${{ steps.ecr.outputs.registry }}
        run: |
          set -euo pipefail
          for app in backend frontend; do
            docker load -i "images/$app.tar"
            docker tag "misiops-$app:latest" "$REGISTRY/misiops-$app:$SHA"
            docker push "$REGISTRY/misiops-$app:$SHA"
            digest=$(docker inspect --format '{{index .RepoDigests 0}}' "$REGISTRY/misiops-$app:$SHA" | cut -d@ -f2)
            echo "${app}_digest=$digest" >> "$GITHUB_OUTPUT"
          done

      - uses: pnpm/action-setup@v4
        with:
          version: 9.6.0
      - uses: actions/setup-node@v4
        with:
          node-version: 20
          cache: pnpm
          cache-dependency-path: infra/pnpm-lock.yaml
      - run: pnpm install --frozen-lockfile
        working-directory: infra

      - name: cdk deploy MisiOps-${{ needs.resolve.outputs.env }}
        working-directory: infra
        run: |
          pnpm cdk deploy "MisiOps-$ENV_NAME" \
            -c "env=$ENV_NAME" -c "imageTag=$SHA" -c "alertEmail=${{ vars.ALERT_EMAIL }}" \
            --require-approval never --outputs-file outputs.json
          echo "PUBLIC_URL=$(jq -r ".\"MisiOps-$ENV_NAME\".PublicUrl" outputs.json)" >> "$GITHUB_ENV"

      - name: Smoke test (FR-017, NFR-009)
        id: smoke
        run: bash infra/scripts/smoke.sh "$PUBLIC_URL" | tee smoke.log

      - name: Summary
        if: always()
        env:
          BACKEND_DIGEST: ${{ steps.push.outputs.backend_digest }}
          FRONTEND_DIGEST: ${{ steps.push.outputs.frontend_digest }}
          SMOKE: ${{ steps.smoke.outcome }}
        run: |
          minutes=$(( ($(date +%s) - START) / 60 ))
          {
            echo "## Deployment to \`$ENV_NAME\`"
            echo "- commit: \`$SHA\` (devsecops run ${{ needs.resolve.outputs.run_id }})"
            echo "- backend digest: \`${BACKEND_DIGEST:-n/a}\`"
            echo "- frontend digest: \`${FRONTEND_DIGEST:-n/a}\`"
            echo "- public URL: ${PUBLIC_URL:-n/a}"
            echo "- smoke test: $SMOKE"
            echo "- duration: ${minutes} min (target: 15 min for updates, 30 min for the first creation)"
          } >> "$GITHUB_STEP_SUMMARY"
\end{lstlisting}

\subsection{Smoke test del despliegue: \texttt{infra/scripts/smoke.sh}}
\label{anx:smoke}

\begin{lstlisting}
#!/usr/bin/env bash
# Post-deployment verification (FR-017, NFR-009, SC-009): the public address serves the web app and the API,
# and both answer within 2 s at p95 over 20 sequential requests. Exit 0 when every check passes, 1 otherwise.
# usage: smoke.sh <public-url>
set -euo pipefail

url=${1:-}
if [ -z "$url" ]; then echo "usage: $0 <public-url>" >&2; exit 2; fi
url=${url%/}
samples=20
p95_limit=2.0
body=$(mktemp); trap 'rm -f "$body"' EXIT

fail() { echo "FAIL - $*"; exit 1; }
# Prints "<status> <seconds>"; a transport error (DNS, TLS, timeout) counts as status 000.
probe() { curl -s -o "$body" -w '%{http_code} %{time_total}' --max-time 10 "$1" || echo "000 10"; }

content_check() { # path expected-substring
  local result status
  result=$(probe "$url$1"); status=${result%% *}
  [ "$status" = "200" ] || fail "GET $1 returned HTTP $status (expected 200)"
  grep -q -- "$2" "$body" || fail "GET $1 body does not contain $2"
  echo "ok   - GET $1 is 200 and contains $2"
}
latency_check() { # path
  local times=() result status
  for _ in $(seq "$samples"); do
    result=$(probe "$url$1"); status=${result%% *}
    [ "$status" = "200" ] || fail "GET $1 returned HTTP $status during sampling"
    times+=("${result#* }")
  done
  [ "${#times[@]}" -eq "$samples" ] || fail "collected ${#times[@]} samples for $1 (expected $samples)"
  local p95
  p95=$(printf '%s\n' "${times[@]}" | sort -g | sed -n "$((samples * 95 / 100))p")
  if awk -v p="$p95" -v l="$p95_limit" 'BEGIN { exit !(p > l) }'; then
    fail "GET $1 p95 ${p95}s over $samples requests exceeds ${p95_limit}s"
  fi
  echo "ok   - GET $1 p95 ${p95}s over $samples requests"
}

content_check / MisiOps
content_check /openapi.json '"openapi"'
latency_check /
latency_check /openapi.json
echo "smoke: all checks passed for $url"
\end{lstlisting}
